% Generated draft: requires mathematical and source review.
\documentclass[11pt]{article}
\usepackage[margin=25mm]{geometry}
\usepackage{fontspec}
\setmainfont{Latin Modern Roman}
\usepackage{amsmath,amssymb,mathtools,unicode-math}
\setmathfont{Latin Modern Math}
\usepackage{enumitem,needspace,xurl,hyperref}
\setlength{\parindent}{0pt}
\setlength{\parskip}{5pt}
\setlength{\emergencystretch}{3em}
\newcommand{\N}{\mathbb{N}}
\newcommand{\Z}{\mathbb{Z}}
\newcommand{\Q}{\mathbb{Q}}
\newcommand{\R}{\mathbb{R}}
\newcommand{\C}{\mathbb{C}}
\newcommand{\F}{\mathbb{F}}
\newcommand{\A}{\mathbb{A}}
\newcommand{\PP}{\mathbb{P}}
\newcommand{\E}{\mathbb{E}}
\newcommand{\eps}{\varepsilon}
\newcommand{\dd}{\,\mathrm{d}}
\begin{document}
\section*{UnsolvedMath ID: 20001857}
\textbf{Source problem code:} catalogue entry
\subsection{Definitions and mathematical statement}
Hempel's source question asks which combined face-angle and dihedral-angle collections are realizable by simplicial polyhedra of fixed combinatorics in \(\mathbb R^3\), and proposes image dimension \(E-1\), explicitly adding ``Needs more clarification.'' Genus, convexity, embeddedness, labels, angle orientation and quotient conventions are not fixed there.

The paired aid chooses a finite connected triangulation \(K\) of a closed orientable genus-\(g\) surface. Let \(V,E,F\) be its numbers of vertices, edges and faces. Consider labeled vertex realizations in \(\mathbb R^3\) with every triangular face nondegenerate, modulo orientation-preserving similarities. Unless separately required, simplexwise-linear realization need not mean globally embedded realization. Record every face angle and each oriented dihedral angle, the latter naturally valued modulo \(2\pi\). Let \(\mathcal A_K\) be this combined-angle map.

The reported reconstruction theorem says that \(\mathcal A_K\) is injective on this quotient and has a local analytic left inverse, using local branches of angle coordinates. Hence, for a nonempty realization model, its actual angle-data image has dimension
\[
 3V-7=E-1-6g,
\]
using \(3F=2E\) and \(V-E+F=2-2g\). The \(E-1\) dimension proposed in the workshop is thus recovered for \(g=0\), but not for positive genus in this model. The aid's example is the embedded Cs\'asz\'ar torus with \(V=7\), \(E=21\): image dimension \(14\), rather than \(20\). Unoriented angles or a different similarity quotient require separate treatment, especially for reflected realizations.

The reported finite realizability criterion develops the triangles along a spanning tree of the dual graph. Begin with a normalized root triangle, with positive face angles summing to \(\pi\), and attach successive triangles using their face angles and prescribed oriented hinge angles. After this development, require equality of coordinates at all repeated occurrences of each labeled vertex, and agreement at every remaining hinge with the prescribed dihedral data. The aid describes these as vertex- and hinge-closure equations, necessary and sufficient in its nondegenerate simplexwise-linear model. Exact equation formulas and analytic charts are not printed in the supplied record. If embeddedness is required, avoidance of unintended intersections must additionally be checked; closure alone does not impose it.

The summary interprets the \(6g\) deficit as global extrinsic closure beyond the \(E-1\)-dimensional space of positive intrinsic edge lengths modulo scale. It does not assert that every such intrinsic metric embeds, nor that the original ambiguous question necessarily intended positive-genus or nonconvex surfaces.
\subsection{Short English statement}
Under the supplied closed orientable labeled-surface model, combined face and oriented dihedral angles reconstruct a realization up to orientation-preserving similarity. The reported image dimension is \(E-1-6g\), recovering \(E-1\) on a sphere and giving a genus-dependent correction otherwise. A dual-tree development gives closure tests; embeddedness and the original intended model still need specification.
\subsection{Sources}
\begin{itemize}
\item Local source record: catalogue entry, catalogue entry; canonical \texttt{aim-geometry-notes.json}, \texttt{notes[194]}, source item 8. Title: ``Angle-data reconstruction and the 6g genus deficit.''
\item AIM 2014 workshop ``Configuration spaces of linkages,'' Problem 8, attributed to Hempel; supplied source reference \url{https://aimath.org/pastworkshops/linkagesproblems.pdf}. This draft uses the supplied record without accessing a PDF. Original remarks and literature fields are empty.
\item Supplied research attempt 1: analytic combined-angle reconstruction, dimension formula, dual-tree closure test and Cs\'asz\'ar torus example. No authoritative external bibliographic entry is supplied for the candidate theorem or example.
\item Registry dated 14 August 2026: ``solved,'' ``SOLVED-BY-YOU,'' statement status ``exact''; research-run status partially solved, result ``counterexample,'' rigor label ``proved,'' novelty confidence low. These are source metadata, not authorship or review claims of this draft. The machine-generated, non-peer-reviewed aid requires independent verification before citation.
\end{itemize}
\end{document}
